\item The discriminant of the degree-five polynomial in \eqref{maps:eq:F.7} is exactly
$\lambda^{15}(3125\lambda+256)$.  At $\lambda_*$ its unique double
root is $r=64/625$ and the coefficient of $(x-r)^2$ is
$10r^3\ne0$.  Its other roots are simple.  Thus the fibre has one
ordinary node and irreducible genus-one normalization.  The local
model is a smoothing of a nonseparating node; the vanishing circle
is a primitive nonseparating homology class.  Cutting and regluing
an annular neighborhood proves that its monodromy is a single
Dehn twist, and hence a primitive rank-one transvection on $H^1$.
In particular its logarithm $N_H$ has square zero and rank one.


\item We identify the representation, rather than merely its spectra.
Choose positive peripheral generators $T_1,T_2$ at infinity and
zero, in that order, with $T_1T_2$ inverse to nodal monodromy.
Set $R=T_1^{-1}$ and $S=T_2$.  Then $R-S$ has rank one and their
characteristic polynomials are $f,g$.  The following elementary
rank-one lemma fixes their simultaneous rational conjugacy class.
If $u$ spans $\operatorname{im}(R-S)$, it is cyclic for $R$:
the span of its iterates is stable under both $R$ and $S$, and
the irreducibility of $f$ makes that span all of $\mathbb Q^4$.
In the rational basis $u,Ru,R^2u,R^3u$, the matrix of $R$ is
the companion matrix of $f$, and $S=R-u\varphi$ for a covector
$\varphi$.  The determinant identity
\[
 g(T)=f(T)+\varphi\operatorname{adj}(TI-R)u
 \tag{F.9}\label{maps:eq:F.9}
\]
determines $\varphi$ uniquely: the four coordinates of
$\operatorname{adj}(TI-R)u$ are a basis of the polynomials of
degree at most three, as is seen from their successively distinct
leading degrees in the cyclic basis.  It follows that any two
such pairs are simultaneously rationally conjugate.  Apply the
lemma to the explicit companion pair
\[
 C_f=\begin{pmatrix}
 0&0&0&-1\\1&0&0&-1\\0&1&0&-1\\0&0&1&-1
 \end{pmatrix},\qquad
 C_g=\begin{pmatrix}
 0&0&0&-1\\1&0&0&0\\0&1&0&0\\0&0&1&0
 \end{pmatrix}.
 \tag{F.10}\label{maps:eq:F.10}
\]
Their difference has rank one and their characteristic polynomials
are $f,g$.  Thus $(T_1,T_2)$ is rationally conjugate to
$(C_f^{-1},C_g)$, with the chosen peripheral conventions.

